\section{Discussion}
\label{sec:discussion}

The oracle experiment (Table~\ref{tab:main_results}, bottom) is not merely an upper bound on SICA performance; it constitutes a causal intervention that isolates the source of failure. By replacing only the constraint source (from self-extracted to gold FOL annotations) while holding all other pipeline components fixed, we establish that the MAX-SAT aggregation mechanism is effective when supplied with discriminative constraints. The +8.82pp improvement ($p < 0.001$) confirms that aggregation is not the bottleneck. This oracle design follows established methodology in compositional reasoning research: \citet{press2023compositionality} use oracle sub-question decompositions to isolate compositionality gaps, \citet{trivedi2023ircot} employ oracle retrieval to separate retrieval from reasoning errors, and \citet{saparov2023greedy} supply gold FOL proofs to test reasoning capacity independent of formalization skill. Our oracle applies the same logic to constraint extraction, pinpointing same-model self-extraction as the single point of failure.

Several recent systems report positive results on overlapping benchmarks using symbolic constraints. DyLoS \citep{xu2026dylos} achieves gains through multi-paradigm solver selection that routes problems to appropriate formal systems. Aristotle \citep{li2025aristotle} applies logic-complete decomposition that breaks complex premises into independently verifiable units. FoVer \citep{wang2025fover} trains a dedicated FOL verifier on supervised data, so the reasoning model need not extract its own constraints. LogicAgent \citep{luo2025logicagent} employs multi-perspective grounding to construct constraints from diverse viewpoints. The common thread is that none of these methods rely on same-model self-extraction. Their success is complementary to our negative finding, not contradictory. The oracle result confirms that the gap these methods bridge (up to +8.82pp) is real and recoverable through external means.

The +8.82pp oracle ceiling quantifies the potential that remains locked behind the capability trap. Three concrete paths can unlock this potential. First, trained FOL extractors, supervised on datasets with gold annotations such as FOLIO, can learn extraction patterns that exceed same-model capabilities. Second, external provers that verify constraints independently of the reasoning model eliminate the circular dependency between generation and verification. Third, cross-model extraction, where a separate model extracts constraints from another model's reasoning traces, breaks the self-reference loop that produces non-discriminative constraints. Each approach targets the extraction bottleneck identified by our ablation: since MAX-SAT and count-only produce identical predictions on 803 of 804 instances (Table~\ref{tab:ablation}), improving the aggregation mechanism offers no path forward.

Our analysis also reveals $C_0$ (the multi-candidate rate) as a zero-cost diagnostic for practitioners considering constraint-based voting. $C_0$ is computable directly from the trace sampling stage, before any constraint extraction occurs. When $C_0$ falls below 5\%, as with QwQ-32B ($C_0 = 2.6\%$), SICA degenerates to SC by construction because nearly all questions receive unanimous answers. High $C_0$ is necessary but not sufficient for SICA to outperform SC, since the capability trap can render extracted constraints non-discriminative even when candidate disagreement is abundant (e.g., LLaMA-3.1-8B on FOLIO, $C_0 = 61.8\%$, $\Delta = 0$). Practitioners should compute $C_0$ from the sampling stage before investing in constraint extraction: low $C_0$ rules out any possible gain, while high $C_0$ warrants further investigation of constraint quality before deployment.
